\bookchapter{Kapitel 69: Spickzettel und Ressourcen}

\chapterintro{Dieses Kapitel ist ein Spickzettel, der auch für sich allein funktioniert. Wenn du dir eine Schreibweise nicht mehr merken kannst oder schnell eine Entscheidung brauchst, blätter einfach direkt hierher. Nutze es als „Taschenkarte“: Code-Spickzettel für den Alltag, häufige Fragen und Lernressourcen, Seite für Seite.}

\booksection{Code-Spickzettel für den Alltag}

\begin{center}
\begin{tabularx}{\linewidth}{|X|X|}
\hline
\textbf{Was du tun willst} & \textbf{Schlüsselcode} \tabularnewline
\hline
Seiteneigener Zustand & @State private var x = startwert \tabularnewline
Seitenübergreifende Logik & \texttt{@Observable class VM \{ var x = 0 \}} \tabularnewline
Layout: untereinander/nebeneinander & \texttt{VStack \{ \}}/\texttt{HStack \{ \}} \tabularnewline
Flexibler Zwischenraum & \texttt{Spacer()} \tabularnewline
Liste & List \{ ForEach(array) \{ ... \} \} \tabularnewline
Seitenwechsel & \texttt{NavigationStack + NavigationLink} \tabularnewline
Unterer Dialog & .sheet(isPresented: \$show) \{ inhalt \} \tabularnewline
Aktuelle Seite schließen & \texttt{@Environment(\textbackslash{}.\allowbreak dismiss) var dismiss} \tabularnewline
Daten aus dem Netz holen & \texttt{try await URLSession.\allowbreak shared.\allowbreak data(from: url)} \tabularnewline
JSON parsen & \texttt{struct X: Codable}+\texttt{JSONDecoder().\allowbreak decode} \tabularnewline
Kleine Einstellungen speichern & UserDefaults.standard.set(wert, forKey: "schlüssel") \tabularnewline
Strukturierte Daten speichern & \texttt{@Model}+\texttt{modelContainer}+\texttt{@Query} \tabularnewline
Dunkelmodus-Anpassung & Systemfarben verwenden (.primary / Color(.systemGray6)), nichts fest im Code verankern \tabularnewline
Tests ausführen & CmdU \tabularnewline
\hline
\end{tabularx}
\end{center}

\begin{center}\small\color{gray}iOS-Entwicklung: Spickzettel (Ausgabe August 2026)\end{center}

\booksection{Häufige Fragen (FAQ)}

F1: Ich habe keinen Mac – kann ich die iOS-Entwicklung trotzdem lernen?

Die Sprache kannst du lernen, eine komplette Entwicklung aber nicht. Mit Swift Playgrounds auf dem iPad/iPhone kannst du den gesamten Stoff von Kapitel 61 durchnehmen; Xcode gibt es jedoch exklusiv für macOS, am Ende brauchst du also einen Mac. Bei knappem Budget kannst du übergangsweise einen Cloud-Mac mieten oder einen gebrauchten Mac mini kaufen und erst nach dem Erlernen der Sprache entscheiden, wie viel du investierst.

F2: Wie viel kostet es, iOS-Entwicklung zu lernen?

Lernphase: 0 Yuan (Xcode ist kostenlos + kostenloses Debuggen auf echten Geräten mit einer kostenlosen Apple-ID). Veröffentlichungsphase: Entwicklerkonto 99 US-Dollar pro Jahr (ca. 700 Yuan). Dazu die Hardware für einen Mac – das sind die gesamten Kosten. Gemessen am Ertrag – selbst ein kleines Tool im App Store hat zahlende Nutzer – geht diese Rechnung auf.

F3: Was lerne ich zuerst – SwiftUI oder UIKit?

SwiftUI. Neue Projekte, neue Tutorials und neue Stellen setzen 2026 auf SwiftUI, und dieses Buch nutzt es durchgehend. UIKit ist das Framework der vorherigen Generation und steckt noch in vielen älteren Projekten – sobald du SwiftUI sicher beherrschst, holst du UIKit bei Bedarf nach (Apple bietet offizielle Interoperabilitäts-Tutorials an), das geht dann recht schnell. Dreh die Reihenfolge nicht um.

F4: Lässt sich mit iOS-Entwicklung gut ein Job finden?

Die Lage 2026: Einstiegsstellen, bei denen es nur darum geht, „Oberflächen zu malen“, sind hart umkämpft – aber Entwickler, die Swift + SwiftUI beherrschen und etwas im Store veröffentlicht haben, sind weiterhin begehrt. Das mittlere Gehalt bei iOS-Stellen liegt in allen Tech-Stacks seit Jahren in der Spitzengruppe. Als Lernender ist dein bester Lebenslauf die App, die du selbst gebaut hast und die im Store steht.

F5: Was tun, wenn die App in der Prüfung abgelehnt wird?

Keine Panik, das ist ganz normal – die Wahrscheinlichkeit, beim ersten Einreichen abgelehnt zu werden, ist durchaus hoch. Apples Ablehnungsgründe sind meist sehr konkret formuliert (welche Seite, welche Richtlinie); verbessere alles Punkt für Punkt, antworte mit einer Erläuterung in App Store Connect und reiche einfach neu ein. Die drei häufigsten Ablehnungsgründe: Der Verwendungszweck von Berechtigungen ist unklar, die Oberfläche hat Bugs, oder die Metadaten (Screenshots/Beschreibung) passen nicht zur App. Wenn du diese drei Punkte vorher selbst prüfst, sparst du dir eine Runde Wartezeit.

\booksection{Die wichtigsten Xcode-Tastenkürzel}

\begin{center}
\begin{tabularx}{\linewidth}{|X|X|}
\hline
\textbf{Aktion} & \textbf{Tastenkürzel} \tabularnewline
\hline
Ausführen / Testen & CmdR / CmdU \tabularnewline
Canvas-Vorschau öffnen & OptionCmd \tabularnewline
Datei schnell öffnen & CmdO \tabularnewline
Globale Suche & CmdF \tabularnewline
Einrückung formatieren & CtrlI \tabularnewline
Kommentieren / Kommentar aufheben & Cmd/ \tabularnewline
Build-Cache leeren & CmdK \tabularnewline
Zur Definition springen & Cmd-Klick \tabularnewline
Umbenennen (sicher, projektweit) & Rechtsklick Refactor Rename \tabularnewline
\hline
\end{tabularx}
\end{center}

\begin{center}\small\color{gray}Die wichtigsten Xcode-Tastenkürzel\end{center}

\begin{tcbinfo}{Tipp: Bei mysteriösen Build-Fehlern zuerst den Cache leeren}
Bei den Spukphänomenen der Art „gestern lief es noch, heute nicht mehr“ liegt es in sechs bis sieben von zehn Fällen an verhunztem Build-Cache. Mit CmdK bereinigen und neu bauen – das heilt die meisten dieser Rätsel. Hilft das nicht, starte Xcode neu (am besten gleich mit dem Mac), und wenn das auch nichts nützt: „Lösch den DerivedData-Ordner“ – der legendäre Dreischlag, der unter iOS-Entwicklern von Mund zu Mund geht.

\end{tcbinfo}

\booksection{Empfohlene Lernressourcen}

\booksection{Empfohlene Lernressourcen}

\begin{center}
\begin{tabularx}{\linewidth}{|X|X|X|}
\hline
\textbf{Typ} & \textbf{Ressource} & \textbf{Empfehlungsgrund} \tabularnewline
\hline
Offiziell & developer.apple.com & Die offizielle Apple-Dokumentation; die SwiftUI-/Swift-Tutorials stecken alle darin – die Autorität Nummer eins \tabularnewline
Offiziell & Apples SwiftUI Tutorials & Der offizielle Schritt-für-Schritt-Kurs, maximale Anfängerfreundlichkeit \tabularnewline
Offiziell & Swift-Website (swift.org) & Offizielle Sprachdokumentation und Online-Playground \tabularnewline
Community & Hacking with Swift & Paul Hudsons kostenlose Tutorial-Sammlung, deren Qualität allgemein als absolute Spitzenklasse gilt \tabularnewline
Community & Stack Overflow (Tag „swiftui“) & Die erste Anlaufstelle bei Fehlermeldungen \tabularnewline
Video & Der offizielle Apple Developer-Kanal (WWDC-Videos) & Die maßgeblichen Erklärungen zu den neuen Technologien des jeweiligen Jahres \tabularnewline
Design & Apples Human Interface Guidelines (HIG) & Die offizielle Bibel für Oberflächen mit echtem „Apple-Flair“ \tabularnewline
Begleitend & Teil 7 „KI-Programmierung in der Praxis“ (Kapitel 37–47) & Die komplette Vorgehensweise, um iOS-Entwicklung mit KI-Unterstützung zu meistern \tabularnewline
\hline
\end{tabularx}
\end{center}

\begin{center}\small\color{gray}Liste der Lernressourcen für die iOS-Entwicklung (August 2026)\end{center}

\booksection{Fazit: vom ersten App zum N-ten App}

Blick zurück: Du bist von „ich weiß nicht mal, was SwiftUI ist“ bis hierher gekommen – Xcode installiert, Swift gelernt, mit SwiftUI Oberflächen gebaut, dem App Datenbank und Netzwerk angeschlossen, und schließlich getestet, gepackt und im Store veröffentlicht. Auf diesem Weg hast du den schwersten Abschnitt von null auf eins bereits gemeistert.

Die Welt der Technik ändert sich schnell: Frameworks werden aktualisiert, Werkzeuge ausgetauscht, und die KI übernimmt immer mehr der „Tastaturarbeit“. Aber eines veraltet nie: deine Fähigkeit, eine vage Idee in klare Schritte zu zerlegen und sie dann Stück für Stück umzusetzen. Werkzeuge sind Verstärker, die Denkweise ist das Fundament. Verstärker kann man austauschen, das Fundament darf nicht wackeln.

Du bist nicht mehr nur "jemand, der mit dem Handy umgehen kann": Du kannst die Ideen in deinem Kopf in Apps verwandeln, die wirklich auf Hunderten Millionen iPhones laufen. Ab heute erschaffst du mit der Tastatur – das ist die Romantik des iOS-Entwicklers.

\booksection{Literaturverzeichnis}

\begin{enumerate}
  \item Apple Inc. (2026). SwiftUI Framework Documentation. developer.apple.com.
  \item Apple Inc. (2026). Swift Programming Language Guide. swift.org.
  \item Apple Inc. (2026). SwiftData Framework Documentation. developer.apple.com.
  \item Apple Inc. (2026). App Store Review Guidelines. developer.apple.com/app-store/review.
  \item Apple Inc. (2026). Human Interface Guidelines (HIG). developer.apple.com/design.
  \item Apple Inc. (2026). Xcode 27 Release Notes and App Distribution Guide.
\end{enumerate}
